\section{提出的方法}

在保持 Tucker 低秩骨架的前提下，如何以受控方式引入高阶交互，并同时保持参数规模、可解释性与优化可实现性？为此，本章先给出 NTD-PL 的结构动机，再正式定义模型，随后概括其基本性质，并给出与当前实现一致的统一学习目标和优化框架。

\subsection{设计动机：显式高阶交互}

标准 Tucker 分解 $\mathcal{T}(\mathcal{X},\{\mathbf{A}^{(n)}\})$ 以共享核心张量和各模因子矩阵构成多线性生成机制，其本质是在低秩潜在空间中通过连续模乘生成观测张量。该结构在压缩、恢复与表示学习中具有良好的参数效率与结构透明性，但它默认观测条目可以由多线性机制充分刻画。对于存在饱和响应、非线性观测、复杂高阶耦合或曲面结构的数据，仅依赖多线性假设往往会带来系统性失配。

在矩阵情形中，一个自然的增强思路是从线性映射
\begin{equation}
\mathbf y = \mathbf A\mathbf x
\end{equation}
推广到带有二次交互的映射
\begin{equation}
\label{eq:qmf_motivation}
\mathbf y = \mathbf A\mathbf x + \mathcal A(\mathbf x,\mathbf x).
\end{equation}
QMF \cite{zhai2024quadratic} 的核心启发正是：二次项的意义不在于单纯增参，而在于它提供了从局部平坦表示走向局部曲面表示的最基本一步。沿着同样的思路，对于张量情形，一个自然问题是：既然 Tucker 已经是线性低秩模型在多模数据上的推广，能否进一步将显式高阶交互引入 Tucker 的生成机制？

为此本文把 $p$ 阶交互写入 Tucker，一个自然的显式原型可写为
\begin{equation}
\label{eq:explicit_p_tucker_prototype}
\widehat y^{(p)}_{i_1\cdots i_N}
=
\sum_{\mathbf r^{(1)},\dots,\mathbf r^{(p)}}
\Bigl(\prod_{q=1}^{p} x_{\mathbf r^{(q)}}\Bigr)
\prod_{n=1}^{N}
 a^{(n)}_{i_n,\,r_n^{(1)},\dots,r_n^{(p)}},
\end{equation}
其中每个模因子在第 $n$ 模上同时依赖 $p$ 个潜在索引。式~\eqref{eq:explicit_p_tucker_prototype} 定义了一种高次 Tucker 分解：它保留了逐模生成的 Tucker 外形，但把每一模的线性映射替换成了显式的 $p$ 阶局部耦合。对于 $p=2$，它对应二次 Tucker；一般 $p$ 时，可视为一个显式 $p$ 次 Tucker 原型，我们记为 $p$-Tucker。下图用张量网络更直观地展示这种 Tucker 的高次扩展。
\begin{center}
\resizebox{\linewidth}{!}{\begin{tikzpicture}[
    line cap=round,
    line join=round,
    >=Latex,
    font=\small,
    factor/.style={
        circle,
        draw=black,
        fill=white,
        minimum size=9mm,
        inner sep=0pt,
        line width=1.0pt
    },
    core/.style={
        circle,
        draw=black,
        fill=gray!20,
        minimum size=9mm,
        inner sep=0pt,
        line width=1.0pt
    },
    edge/.style={
        draw=black,
        line width=1.0pt
    },
    hedge/.style={
        draw=black,
        dashed,
        line width=1.0pt,
        dash pattern=on 4pt off 3pt
    }
]

% =========================================================
% Tucker  (3D-like view)
% =========================================================
\begin{scope}[xshift=0cm]
    % outer factor nodes
    \node[factor] (f1) at (0.00, 0.00) {$\mathbf{A}^{(1)}$};
    \node[factor] (f2) at (3.464, 2.000) {$\mathbf{A}^{(2)}$};
    \node[factor] (f3) at (2.249, 2.932) {$\mathbf{A}^{(3)}$};

    % core node
    \node[core]   (c1) at (1.732, 1.00) {$\mathcal{X}$};

    % edges
    \draw[edge] (f1) -- (c1);
    \draw[edge] (f2) -- (c1);
    \draw[edge] (f3) -- (c1);

    \node[font=\normalsize] at (1.732,-1.20) {Tucker};
\end{scope}

\node[font=\Large] at (4.00, 1.00) {$\Longrightarrow$};

% =========================================================
% Quadratic Tucker
% =========================================================
\begin{scope}[xshift=5cm]
    % outer factor nodes
    \node[factor] (f1) at (0.00, 0.00) {$\mathcal{A}^{(1)}$};
    \node[factor] (f2) at (3.464, 2.000) {$\mathcal{A}^{(2)}$};
    \node[factor] (f3) at (2.249, 2.932) {$\mathcal{A}^{(3)}$};

    % inner/core nodes
    \node[core] (c1) at (1.378, 1.354) {$\mathcal{X}$};
    \node[core] (c2) at (2.085, 0.64) {$\mathcal{X}$};

    % visible structure
    \draw[edge] (f1) -- (c1);
    \draw[edge] (f1) -- (c2);
    \draw[hedge] (f2) -- (c1);
    \draw[edge] (f2) -- (c2);
    \draw[edge] (f3) -- (c1);
    \draw[edge] (f3) -- (c2);

    \node[font=\normalsize] at (1.732,-1.20) {Quadratic Tucker};
\end{scope}

\node[font=\Large] at (9.00, 1.00) {$\Longrightarrow$};

% =========================================================
% p-Tucker
% =========================================================
\begin{scope}[xshift=10cm]
    % outer factor nodes
    \node[factor] (f1) at (0.00, 0.00) {$\mathcal{A}^{(1)}$};
    \node[factor] (f2) at (3.464, 2.000) {$\mathcal{A}^{(2)}$};
    \node[factor] (f3) at (2.249, 2.932) {$\mathcal{A}^{(3)}$};

    % inner/core nodes
    \node[core] (c1) at (1.378, 1.354) {$\mathcal{X}$};
    \node[core] (c3) at (0.671, 2.061) {$\mathcal{X}$};
    \node[core] (c4) at (2.793, -0.061) {$\mathcal{X}$};

    % main visible edges
    \draw[edge] (f1) -- (c1);
    \draw[edge] (f1) -- (c3);
    \draw[edge] (f1) -- (c4);
    \draw[hedge] (f2) -- (c1);
    \draw[hedge] (f2) -- (c3);
    \draw[edge] (f2) -- (c4);
    \draw[edge] (f3) -- (c1);
    \draw[edge] (f3) -- (c3);
    \draw[edge] (f3) -- (c4);

    % vertical continuation hint
    \node[font=\Large, rotate=-45] at (2.085, 0.64) {$\dots$};

    \node[font=\normalsize] at (1.732,-1.20) {$p$-Tucker};
\end{scope}

\end{tikzpicture}}
\end{center}

这一原型虽然在概念上自然，却不适合作为最终模型。随着次数 $p$ 增大，模因子需要提升为更高阶张量，核心张量在生成过程中也会诱导出 Kronecker 型高阶扩展，参数量、存储与收缩计算都会迅速膨胀。更重要的是，优化难度也会随之显著增加。因此，显式 $p$-Tucker 更适合作为结构动机：它解释了高阶交互进入 Tucker 的方式，却没有直接给出一个可扩展的实现。

基于这一观察，本文不直接训练式~\eqref{eq:explicit_p_tucker_prototype} 所对应的显式高阶模型，而是寻找一种结构化压缩：既保留 $p$-Tucker 高阶交互的语义，又保持 Tucker 的参数量。NTD-PL 正是这一设计目标下得到的模型。

\subsection{NTD-PL 模型定义与结构解释}

\subsubsection{模型定义}

本文提出多项式链接的非线性 Tucker 分解 NTD-PL。其核心思想是：先用 Tucker 分解生成一个共享的低秩潜在张量，再在该潜在张量上施加显式的逐元素多项式链接。具体而言，记
\begin{equation}
\label{eq:ntdpl_latent_tensor}
\cS
=
\mathcal{T}(\mathcal{X},\{\mathbf{A}^{(n)}\}),
\end{equation}
其中 $\cX\in\mathbb R^{R_1\times\cdots\times R_N}$ 为核心张量，$\mathbf A^{(n)}\in\mathbb R^{I_n\times R_n}$ 为第 $n$ 模因子矩阵。NTD-PL 定义为
\begin{equation}
\label{eq:ntdpl_model}
\widehat{\cY}
=
 f_{\boldsymbol\beta}(\cS)
=
\sum_{p=0}^{P} \beta_p \cS^{\odot p},
\end{equation}
其中 $\boldsymbol\beta=(\beta_0,\beta_1,\dots,\beta_P)^\top$ 为多项式链接系数，$\cS^{\odot p}$ 表示 $\cS$ 的逐元素 $p$ 次幂。

式~\eqref{eq:ntdpl_model} 中，$\beta_0$ 对应常数偏置项，$\beta_1$ 对应标准 Tucker 线性项，而 $\beta_p$（$p\ge 2$）则对应在共享低秩 骨架 上激活的高阶交互项。特别地，当
\begin{equation}
\beta_1=1,\qquad \beta_p=0\ (p\neq 1)
\end{equation}
时，模型精确退化为标准 Tucker。因此，NTD-PL 不是替代 Tucker 的另一类新骨架，而是以 Tucker 潜在张量为中心的受控非线性扩展。

NTD-PL 不能简单理解为先做一次 Tucker 重构，再对输出做逐元素校正。关键区别在于：式~\eqref{eq:ntdpl_model} 中的非线性作用对象不是一个与低秩结构无关的外部信号，而是共享 Tucker 潜在张量 $\cS$ 本身。由于 $\cS$ 的每个条目已经由核心张量与各模因子的多线性耦合生成，因此 $\cS^{\odot p}$ 对应的是同一 骨架 上交互模式的重组，而不是对一个已经完成的线性重构施加独立校准。后文将专门通过 Tucker 的多项式后处理 与 NTD-PL 的对比，区分这两种机制。

\subsubsection{作为 $p$-Tucker 的结构化压缩}

NTD-PL 的定义来源于对 $p$-Tucker的一种结构化压缩。以二次情形为例，若将式~\eqref{eq:explicit_p_tucker_prototype} 中的模因子写成
\begin{equation}
 a^{(n)}_{i_n,r_n,r_n'} \approx b^{(n)}_{i_n,r_n}c^{(n)}_{i_n,r_n'},
\end{equation}
则二次项可重写为两个 Tucker 张量的逐元素乘积：
\begin{equation}
\label{eq:separable_quadratic_factorization}
\widehat{\cY}_2 \approx \cS_B \odot \cS_C,
\end{equation}
其中
\begin{equation}
\mathcal{T}(\mathcal{X},\{\mathbf{B}^{(n)}\}),
\qquad
\mathcal{T}(\mathcal{X},\{\mathbf{C}^{(n)}\}).
\end{equation}
进一步施加共享投影约束 $\mathbf B^{(n)}=\mathbf C^{(n)}$，便得到
\begin{equation}
\label{eq:shared_quadratic_term}
\widehat{\cY}_2 \approx \cS^{\odot 2}.
\end{equation}
按同样思路推广到更高阶，就得到共享 骨架 上的 $\cS^{\odot p}$。于是，NTD-PL 可被理解为将这些共享投影的 $p$ 阶交互项按系数线性组合，从而把显式 $p$-Tucker 的复杂性压缩到少量链接系数中，下图用张量网络图直观地表示了这一行为。
\begin{center}
\resizebox{\linewidth}{!}{\begin{tikzpicture}[
    line cap=round,
    line join=round,
    >=Latex,
    font=\small,
    factor/.style={
        circle,
        draw=black,
        fill=white,
        minimum size=9mm,
        inner sep=0pt,
        line width=1.0pt
    },
    core/.style={
        circle,
        draw=black,
        fill=gray!20,
        minimum size=9mm,
        inner sep=0pt,
        line width=1.0pt
    },
    edge/.style={
        draw=black,
        line width=1.0pt
    },
    hedge/.style={
        draw=black,
        dashed,
        line width=1.0pt,
        dash pattern=on 4pt off 3pt
    }
]

% =========================================================
% p-Tucker
% =========================================================
\begin{scope}[xshift=0cm]
    \node[factor] (f1) at (0.00, 0.00) {$\mathcal{A}^{(1)}$};
    \node[factor] (f2) at (3.464, 2.000) {$\mathcal{A}^{(2)}$};
    \node[factor] (f3) at (2.249, 2.932) {$\mathcal{A}^{(3)}$};

    \node[core] (c1) at (1.378, 1.354) {$\mathcal{X}$};
    \node[core] (c2) at (0.671, 2.061) {$\mathcal{X}$};
    \node[core] (c3) at (2.793, -0.061) {$\mathcal{X}$};

    \draw[edge] (f1) -- (c1);
    \draw[edge] (f1) -- (c2);
    \draw[edge] (f1) -- (c3);
    \draw[hedge] (f2) -- (c1);
    \draw[hedge] (f2) -- (c2);
    \draw[edge] (f2) -- (c3);
    \draw[edge] (f3) -- (c1);
    \draw[edge] (f3) -- (c2);
    \draw[edge] (f3) -- (c3);

    \node[font=\Large, rotate=-45] at (2.085, 0.64) {$\dots$};
    \node[font=\normalsize] at (1.732,-1.35) {$p$-Tucker};
\end{scope}

\node[font=\Large] at (4.65, 1.05) {$\approx$};

% =========================================================
% Tucker 1
% =========================================================
\begin{scope}[xshift=5.8cm]
    \node[factor] (g1f1) at (0.00, 0.00) {$\mathbf{B}^{(1)}$};
    \node[factor] (g1f2) at (3.464, 2.000) {$\mathbf{B}^{(2)}$};
    \node[factor] (g1f3) at (2.249, 2.932) {$\mathbf{B}^{(3)}$};
    \node[core]   (g1c1) at (1.732, 1.00) {$\mathcal{X}$};

    \draw[edge] (g1f1) -- (g1c1);
    \draw[edge] (g1f2) -- (g1c1);
    \draw[edge] (g1f3) -- (g1c1);

    \node[font=\normalsize] at (1.732,-1.35) {$\mathcal{S}_1$};
\end{scope}

\node[font=\Large] at (10.05, 1.05) {$\odot$};
\node[font=\Large] at (11.25, 1.05) {$\cdots$};
\node[font=\Large] at (12.45, 1.05) {$\odot$};

% =========================================================
% Tucker p
% =========================================================
\begin{scope}[xshift=13.6cm]
    \node[factor] (gpf1) at (0.00, 0.00) {$\mathbf{P}^{(1)}$};
    \node[factor] (gpf2) at (3.464, 2.000) {$\mathbf{P}^{(2)}$};
    \node[factor] (gpf3) at (2.249, 2.932) {$\mathbf{P}^{(3)}$};
    \node[core]   (gpc1) at (1.732, 1.00) {$\mathcal{X}$};

    \draw[edge] (gpf1) -- (gpc1);
    \draw[edge] (gpf2) -- (gpc1);
    \draw[edge] (gpf3) -- (gpc1);

    \node[font=\normalsize] at (1.732,-1.35) {$\mathcal{S}_p$};
\end{scope}

\end{tikzpicture}
}
\end{center}

\subsection{NTD-PL 的基本性质}

本节概括 NTD-PL 的三个最基本性质：模型族关系与参数效率、非线性生成下的表达能力，以及尺度不定性与归一化机制。详细证明见附录~\ref{app:model_properties}--\ref{app:expressivity}。之后我们会看到，这些性质在可控合成数据实验中得到验证。

\subsubsection{模型族关系与参数效率}

\begin{proposition}[退化到 Tucker]
\label{prop:tucker_special_case}
若 $\beta_1=1$，且对所有 $p\neq 1$ 有 $\beta_p=0$，则式~\eqref{eq:ntdpl_model} 退化为标准 Tucker 模型，即
\begin{equation}
\widehat{\cY}
=
\cS
=
\mathcal{T}(\mathcal{X},\{\mathbf{A}^{(n)}\}).
\end{equation}
\end{proposition}

\begin{proposition}[次数嵌套性]
\label{prop:nested_model_family}
对任意 $P_1<P_2$，degree-$P_1$ 的 NTD-PL 模型类是 degree-$P_2$ 模型类的子集。
\end{proposition}

这两个命题说明，NTD-PL 构成了一族关于多项式次数 $P$ 的嵌套模型：当高阶项关闭时，它精确回到 Tucker；随着 $P$ 增加，模型表达力按次数受控扩张。这也是后文采用逐次延续优化的理论基础。

\begin{proposition}[参数复杂度]
\label{prop:param_complexity}
设多线性秩为 $(R_1,\dots,R_N)$，则 NTD-PL 的参数总数为
\begin{equation}
\#\mathrm{param}
=
\prod_{n=1}^{N}R_n
+
\sum_{n=1}^{N}I_nR_n
+
(P+1),
\label{eq:ntdpl_param_count_main}
\end{equation}
其中前两项对应标准 Tucker 分解的核心与因子参数量，最后一项对应多项式链接系数。
\end{proposition}

命题~\ref{prop:param_complexity} 表明，NTD-PL 的额外自由度只体现为 $P+1$ 个标量系数，而不是新的高阶因子张量或高阶核心。因而，相比显式 $p$-Tucker，NTD-PL 不是通过增加另一套高阶参数化来增强表达力，而是在基本保持 Tucker 参数规模一致的前提下引入受控非线性。

\subsubsection{非线性生成下的表达能力}

NTD-PL 的表达能力可以从逐元素非线性生成模型的角度理解。设存在某个真实 Tucker 潜在张量
\begin{equation}
\cS^\star
=
\mathcal{T}(\mathcal{X^\star},\{\mathbf{A}^{(n)\star}\})
\end{equation}
并设观测张量由某个标量函数 $g:\mathbb R\to\mathbb R$ 逐元素作用生成，即
\begin{equation}
\cY = g(\cS^\star).
\end{equation}

\begin{proposition}[多项式生成下的精确表达]
\label{prop:poly_exact_main}
若观测张量满足 $\cY = g(\cS^\star)$，其中 $g(s)=\sum_{p=0}^{P} a_p s^p$ 为次数不超过 $P$ 的多项式，则取潜在张量 $\cS=\cS^\star$ 且 $\beta_p=a_p$ 时，NTD-PL 可精确表示该生成过程。详细证明见附录定理~\ref{thm:poly_exact}。
\end{proposition}

\begin{theorem}[解析函数下的显式截断误差界]
\label{thm:analytic_bound_main}
设 $g$ 在复圆盘
\begin{equation}
D(0,\rho)=\{z\in\mathbb C:|z|<\rho\}
\end{equation}
内解析，且存在常数 $B$ 满足
\begin{equation}
\|\cS^\star\|_\infty \le B < \rho.
\end{equation}
记 $g$ 在 $0$ 点的 Taylor 展开为
\begin{equation}
g(s)=\sum_{p=0}^{\infty} a_p s^p,
\end{equation}
并记其 $P$ 阶截断多项式为
\begin{equation}
q_P(s)=\sum_{p=0}^{P} a_p s^p.
\end{equation}
若取 NTD-PL 的潜在张量为 $\cS^\star$，并令链接系数满足 $\beta_p=a_p$，则其逼近误差满足
\begin{equation}
\|g(\cS^\star)-q_P(\cS^\star)\|_F
\le
\sqrt{N_e}\,M_\rho
\frac{(B/\rho)^{P+1}}{1-B/\rho},
\label{eq:analytic_bound_main}
\end{equation}
其中
\begin{equation}
N_e=\prod_{n=1}^N I_n,
\qquad
M_\rho=\sup_{|z|=\rho}|g(z)|.
\end{equation}
\end{theorem}

式~\eqref{eq:analytic_bound_main} 给出了一个具有明确解释性的误差分解。首先，主导项 $(B/\rho)^{P+1}$ 表明，当 $B/\rho<1$ 固定时，截断误差会随多项式阶数 $P$ 以几何速度衰减，因此只要真实潜在张量的幅值范围没有逼近解析半径边界，有限阶 NTD-PL 就能够快速逼近解析生成函数。其次，分母项 $1-B/\rho$ 刻画了“距离奇点的裕量”：当 $B$ 相对 $\rho$ 更小时，说明真实信号活动区间位于更稳定的解析区域内部，此时余项界更紧；反之，当 $B$ 逐渐接近 $\rho$ 时，逼近难度会显著上升。最后，$\sqrt{N_e}$ 只反映从标量一致逼近推广到整张量 Frobenius 误差时的规模放大，而不改变误差关于 $P$ 的收敛阶。因而，该定理说明 NTD-PL 的逼近能力主要由链接函数在有效输入区间上的解析性以及 $B/\rho$ 的相对大小决定，而不是由额外高阶张量参数的引入所驱动。这些将会解释第 \ref{sec:nonlinearapprox} 节中的逼近的现象。

\subsubsection{尺度不定性与归一化机制}

与标准 Tucker 一样，NTD-PL 继承了核心张量与因子矩阵之间的尺度不定性。由于 NTD-PL 同时组合 $\cS,\cS^{\odot 2},\dots,\cS^{\odot P}$，这一尺度问题在多项式链接中更加敏感：若 $\cS$ 的尺度漂移，不同阶次项会按不同速率放大或缩小，进而直接影响 $\boldsymbol\beta$ 的数值稳定性。

\begin{proposition}[尺度不变性]
\label{prop:scale_invariance_main}
设对某一模 $n$ 的因子矩阵施加可逆对角缩放 $D^{(n)}$，并同时将对应逆缩放吸收到核心张量中，则潜在张量 $\cS$ 保持不变，从而 NTD-PL 的输出 $\widehat{\cY}=f_{\boldsymbol\beta}(\cS)$ 也保持不变。
\end{proposition}

命题~\ref{prop:scale_invariance_main} 说明，因子列归一化并将尺度吸收到核心张量中，可以理解为一种 \emph{gauge fixing}：它不改变模型能力，只从等价参数化中选取更稳定的表示。在高阶项参与建模时，这一步对 $\boldsymbol\beta$ 的估计尤为重要，因此后文优化算法中将其作为固定步骤。

\subsection{统一学习目标与优化算法}

\subsubsection{统一学习目标}

为在同一优化框架下同时处理分解与缺失补全，记二值观测掩码张量为 $\Omega\in\{0,1\}^{I_1\times\cdots\times I_N}$。考虑如下目标函数：
\begin{equation}
\label{eq:ntdpl_masked_objective}
\min_{\cX,\{\mathbf A^{(n)}\},\boldsymbol\beta}
\;
\frac{1}{2Z_\Omega}
\left\|
\Omega \odot
\left(
\cY - f_{\boldsymbol\beta}(\cS)
\right)
\right\|_F^2
+
\mathcal R(\cX,\{\mathbf A^{(n)}\},\boldsymbol\beta),
\end{equation}
其中 $\cS=\mathcal{T}(\mathcal{X},\{\mathbf{A}^{(n)}\})$，且
\begin{equation}
Z_\Omega =
\begin{cases}
\prod_{n=1}^{N} I_n, & \Omega \equiv \mathbf 1,\\[2mm]
\sum\limits_{i_1,\dots,i_N} \Omega_{i_1\cdots i_N}, & \text{否则}.
\end{cases}
\end{equation}
这里的 $Z_\Omega$ 用于按有效条目数归一化数据拟合项，使完整观测与缺失观测两种情形下的梯度量级保持可比；这也与当前实现中的残差缩放方式一致。正则项取为
\begin{equation}
\label{eq:ntdpl_reg_obj}
\mathcal R(\cX,\{\mathbf A^{(n)}\},\boldsymbol\beta)
=
\frac{\lambda_X}{2}\|\cX\|_F^2
+
\sum_{n=1}^{N}\frac{\lambda_n}{2}\|\mathbf A^{(n)}\|_F^2
+
\frac{\lambda_\beta}{2}\|\boldsymbol\beta\|_2^2.
\end{equation}

当 $\Omega\equiv \mathbf 1$ 时，式~\eqref{eq:ntdpl_masked_objective} 就是完整张量上的 NTD-PL 分解；当 $\Omega$ 为一般掩码时，它对应仅在观测位置上训练的补全目标。由此，分解与补全不再需要两套不同模型，只是同一学习目标下的不同观测设置。

\subsubsection{初始化}

NTD-PL 是 Tucker 的多项式推广，自然可以利用Tucker分解作为优化起点。对完整观测使用标准 Tucker 初始化，对缺失观测使用 masked Tucker 初始化 \cite{yang2016iterative, liu2012tensor}。对多项式系数，采用线性 warm start：
\begin{equation}
\beta_0=0,\qquad \beta_1=1,\qquad \beta_p=0\quad (p\ge 2).
\label{eq:beta_warm_start}
\end{equation}
本文采用逐次延续的优化方式，从 $p=1$ 的线性情形开始，随着优化推进再逐步激活更高阶项，意味着模型总是从标准 Tucker 解附近出发，而不是从一个高阶非线性模型直接起步。

\subsubsection{核心与因子的更新}

固定 $\boldsymbol\beta$ 后，首先构造潜在张量
\begin{equation}
\cS = \llbracket \cX;\mathbf A^{(1)},\dots,\mathbf A^{(N)}\rrbracket,
\end{equation}
并计算预测张量
\begin{equation}
\widehat{\cY}=f_{\boldsymbol\beta}(\cS),
\qquad
f_{\boldsymbol\beta}'(\cS)=\sum_{p=1}^{P} p\beta_p\cS^{\odot(p-1)}.
\label{eq:poly_derivative_tensor}
\end{equation}
在当前实现中，$f_{\boldsymbol\beta}(\cS)$ 与 $f_{\boldsymbol\beta}'(\cS)$ 通过一次融合的 Horner 递推同时求值，以减少重复的幂次计算与内存访问；这只改变常数开销，不改变模型本身。

记归一化后的掩码残差为
\begin{equation}
\mathbf E = \frac{1}{Z_\Omega}\,\Omega\odot(\widehat{\cY}-\cY),
\end{equation}
并记
\begin{equation}
\mathbf T = \mathbf E \odot f_{\boldsymbol\beta}'(\cS).
\label{eq:T_tensor}
\end{equation}
则根据链式法则，目标函数对潜在张量 $\cS$ 的梯度正比于 $\mathbf T$。进一步通过 Tucker 重构关系反向传播，可得核心张量梯度
\begin{equation}
\label{eq:grad_core_main}
\nabla_{\cX}\mathcal L
=
\mathbf T
\times_1 \mathbf A^{(1)\top}
\times_2 \cdots
\times_N \mathbf A^{(N)\top}
+
\lambda_X\cX.
\end{equation}
对于第 $n$ 模因子矩阵，记
\begin{equation}
\label{eq:Zn_def}
\mathbf Z^{(n)}
=
\Bigl[\cX\times_{m\neq n}\mathbf A^{(m)}\Bigr]_{(n)},
\end{equation}
其中 $[\cdot]_{(n)}$ 表示模 $n$ 展开，则相应梯度可写为
\begin{equation}
\label{eq:grad_factor_main}
\nabla_{\mathbf A^{(n)}}\mathcal L
=
\mathbf T_{(n)}\,\mathbf Z^{(n)\top}
+
\lambda_n\mathbf A^{(n)}.
\end{equation}

由于逐元素多项式链接打破了经典 Tucker 中最小二乘子问题的闭式结构，本文不采用交替最小二乘法，而是使用 Adam \cite{kingma2014adam} 对 $\cX$ 与各模因子做梯度更新。

\subsubsection{多项式系数的更新}

固定 $\cX$ 与 $\{\mathbf A^{(n)}\}_{n=1}^{N}$ 时，潜在张量 $\cS$ 也随之固定，此时 $\boldsymbol\beta$ 的更新可归结为一个低维 ridge 回归问题 \cite{hoerl1970ridge}。记观测向量为 $\mathbf y_\Omega$，并定义设计矩阵
\begin{equation}
\label{eq:PhiOmega_main}
\Phi_\Omega
=
\left[
\mathbf 1,
\mathrm{vec}(\cS)_\Omega,
\mathrm{vec}(\cS^{\odot 2})_\Omega,
\dots,
\mathrm{vec}(\cS^{\odot P})_\Omega
\right]
\in \mathbb R^{|\Omega|\times(P+1)}.
\end{equation}
则 $\boldsymbol\beta$ 子问题为
\begin{equation}
\label{eq:beta_ridge}
\min_{\boldsymbol\beta}
\frac12\|\mathbf y_\Omega-\Phi_\Omega\boldsymbol\beta\|_2^2
+
\frac{\lambda_\beta}{2}\|\boldsymbol\beta\|_2^2.
\end{equation}
若允许常数项参与拟合，则直接求解式~\eqref{eq:beta_ridge}；若不允许常数项，则固定 $\beta_0=0$，仅对 $p\ge 1$ 的系数做相同的 ridge 更新。由于 $\lambda_\beta>0$，系统矩阵 $\Phi_\Omega^\top\Phi_\Omega+\lambda_\beta I$ 始终正定，因此该子问题强凸且存在唯一解。其严格证明见附录~\ref{app:beta_unique}。

尽管整体问题对 $(\cX,\{\mathbf A^{(n)}\},\boldsymbol\beta)$ 联合来看是非凸的，但在固定 Tucker 潜变量 $\cS$ 后，对 $\boldsymbol\beta$ 的更新始终是一个规模仅为 $P+1$ 的强凸 ridge 回归。这也是本文能够在保持高阶非线性表达的同时，仍使用稳定交替优化框架的关键原因。

高阶幂基在数值上可能病态，尤其当 $\cS$ 的取值范围较大时更为明显。为改善这一问题，优化实现中可先在标准化变量
\begin{equation}
\mathbf Z = \frac{\cS-\mu}{\sigma}
\end{equation}
的幂基上求解系数 $\boldsymbol\gamma$，再通过基变换将其映回原始幂基系数 $\boldsymbol\beta$。这一过程只改变子问题求解的数值条件，并不改变外层模型的形式。

\subsubsection{因子归一化与核心吸收}

在每轮梯度更新后，本文对每个因子矩阵的列向量做单位 $\ell_2$ 范数归一化，并将对应尺度吸收到核心张量中。根据命题~\ref{prop:scale_invariance_main}，该步骤不会改变潜在张量 $\cS$，因此也不会改变 NTD-PL 的输出。其作用在于减少等价参数化引入的尺度漂移，特别是降低 $\cS$ 的尺度变化对高阶项 $\cS^{\odot p}$ 和 $\boldsymbol\beta$ 估计带来的耦合影响。

\subsubsection{次数延续策略}

直接从零开始优化高阶多项式模型通常不够稳定。为此，本文采用次数延续（degree continuation）策略：训练从 $p=1$ 的线性 Tucker 情形开始，并在预设里程碑 $\tau_2<\cdots<\tau_P$ 处逐步激活更高阶项。每当第 $p$ 阶项被激活时，对应的新系数初始化为零，即
\begin{equation}
\boldsymbol\beta^{\mathrm{new}}
=
(\beta_0,\dots,\beta_{p-1},0)^\top.
\end{equation}
这样，模型便可从稳定的线性解逐步过渡到更丰富的非线性解，而不是一开始就在高阶搜索空间中联合优化全部系数。命题~\ref{prop:nested_model_family} 已说明，这样的 continuation 不是经验性“技巧”，而是顺应模型族嵌套结构的自然优化路径。

\begin{algorithm}[t]
\caption{基于交替优化与次数延续的 NTD-PL}
\label{alg:ntdpl}
\begin{algorithmic}[1]
\State 输入：观测张量 $\cY$、掩码 $\Omega$、多线性秩 $\{R_n\}$、最大多项式次数 $P$
\State 通过 Tucker / masked Tucker 初始化 $(\cX,\{\mathbf A^{(n)}\})$
\State 初始化 $\beta_0=0,\beta_1=1,\beta_p=0\ (p\ge 2)$
\State 若启用 continuation，则当前次数置为 $p\leftarrow 1$；否则置为 $p\leftarrow P$
\For{$t=1,2,\dots,T$}
    \State 若到达新的次数里程碑，则激活下一阶项并将新系数置零
    \State 构造潜在张量 $\cS=\llbracket \cX;\mathbf A^{(1)},\dots,\mathbf A^{(N)}\rrbracket$
    \State 计算 $f_{\boldsymbol\beta}(\cS)$ 与 $f'_{\boldsymbol\beta}(\cS)$，得到残差张量与中间量 $\mathbf T$
    \State 用 Adam 更新核心张量 $\cX$ 与各模因子 $\{\mathbf A^{(n)}\}$
    \State 对因子列归一化，并将尺度吸收到核心张量中
    \If{$t$ 满足 $\boldsymbol\beta$ 更新间隔}
        \State 固定当前 $\cS$，通过 ridge 回归更新 $\boldsymbol\beta$
    \EndIf
\EndFor
\State 输出：$\cX$、$\{\mathbf A^{(n)}\}$ 与 $\boldsymbol\beta$
\end{algorithmic}
\end{algorithm}

\subsubsection{算法复杂度分析}

记观测张量总元素个数为
\begin{equation}
M = \prod_{n=1}^{N} I_n,
\qquad
M_\Omega = |\Omega|,
\qquad
R=\prod_{n=1}^{N}R_n.
\end{equation}
这里 $M$ 表示完整张量的条目总数，$M_\Omega$ 表示观测条目数。

做一次标准 Tucker 重构时，若按逐模乘法顺序实现，其前向复杂度可写为
\begin{equation}
\mathcal C_{\mathrm{rec}}
=
\sum_{n=1}^{N}
I_nR_n
\left(\prod_{m<n} I_m\right)
\left(\prod_{m>n} R_m\right).
\label{eq:tucker_reconstruction_complexity}
\end{equation}
该项给出潜在张量 $\cS$ 的主要构造代价，也是后续梯度传播的基础。

给定 $\cS$ 后，计算 $f_{\boldsymbol\beta}(\cS)$ 与 $f'_{\boldsymbol\beta}(\cS)$ 的代价均为逐元素多项式求值。若采用融合的 Horner 递推，则它们可在一次遍历中同时完成，渐近复杂度为
\begin{equation}
\mathcal O(PM).
\end{equation}
相较于分别计算前向值与导数，这一融合实现仅降低常数因子，不改变关于 $P$ 与 $M$ 的阶数。

核心梯度的计算可视为把张量 $\mathbf T$ 重新投影回核心空间，其复杂度与一次 Tucker 压缩同阶。因子梯度的计算则具有 MTTKRP 型结构：对每个模 $n$，需要先构造式~\eqref{eq:Zn_def} 中的 $\mathbf Z^{(n)}$，再计算 $\mathbf T_{(n)}\mathbf Z^{(n)\top}$。因此，若每轮外迭代对 $(\cX,\{\mathbf A^{(n)}\})$ 执行 $K$ 步 Adam 更新，则其主导成本可近似写为
\begin{equation}
\mathcal O\!\left(
K\bigl(\mathcal C_{\mathrm{rec}} + PM + \mathcal C_{\mathrm{bwd}}\bigr)
\right),
\label{eq:grad_update_complexity}
\end{equation}
其中 $\mathcal C_{\mathrm{bwd}}$ 代表由核心反投影与各模因子 MTTKRP 型乘法构成的反向传播成本，通常与 Tucker 骨架本身的收缩/展开复杂度同阶。

在固定 $\cS$ 时，$\boldsymbol\beta$ 子问题只依赖观测条目。构造观测设计矩阵 $\Phi_\Omega\in\mathbb R^{M_\Omega\times(P+1)}$ 的代价为
\begin{equation}
\mathcal O(PM_\Omega),
\end{equation}
求解正规方程
\begin{equation}
(\Phi_\Omega^\top\Phi_\Omega+\lambda_\beta I)\boldsymbol\beta
=
\Phi_\Omega^\top\mathbf y_\Omega
\end{equation}
的代价为
\begin{equation}
\mathcal O(M_\Omega P^2 + P^3).
\label{eq:beta_update_complexity}
\end{equation}
由于 $P$ 在本文设置中通常很小（如 $2,3,4$），因此 $\boldsymbol\beta$ 更新虽然理论上是额外步骤，但在总体时间中通常并非主导项。

因子归一化与核心吸收的代价为
\begin{equation}
\mathcal O\!\left(\sum_{n=1}^{N} I_nR_n + R\right),
\end{equation}
相较于式~\eqref{eq:grad_update_complexity} 通常可以忽略。次数延续不会改变单轮迭代的渐近复杂度，只会通过逐步激活高阶项影响总迭代轮数。优化实现中，避免重复 Tucker 重构、融合多项式/导数求值、以及在标准化幂基上稳定求解 $\boldsymbol\beta$，都只降低常数开销或改善数值条件，不改变上述渐近阶数。

综合上述各项，当前稠密实现下 NTD-PL 的单轮外迭代复杂度可概括为
\begin{equation}
\mathcal O\!\left(
K\bigl(\mathcal C_{\mathrm{rec}} + PM + \mathcal C_{\mathrm{bwd}}\bigr)
+
PM_\Omega + M_\Omega P^2 + P^3
\right).
\label{eq:total_complexity_main}
\end{equation}
因此，NTD-PL 相比标准 Tucker 的额外计算主要来自两部分：其一是逐元素多项式链接及其导数的计算，其代价为低阶、可并行的 $\mathcal O(PM)$；其二是 $\boldsymbol\beta$ 的小规模 ridge 更新，其代价为 $\mathcal O(PM_\Omega + M_\Omega P^2 + P^3)$。换言之，当前实现下的主导成本仍然来自 Tucker 骨架本身的前向/反向传播，而不是多项式链接系数的求解。
